% TOP_PROBLEM: {"id": 5100002, "title": "Elliptic-billiard invariant k_{108}", "category_id": 11, "category": {"id": 11, "name": "dynamical_systems", "display_name": "Dynamical Systems", "description": "Problems about long-term behavior of deterministic systems, Hamiltonian dynamics, and periodic orbits.", "slug": "dynamical-systems", "order_index": 11, "created_at": "2026-07-31T15:26:25.671Z"}, "status": "open"}
% FIRST_POSTED: null
% ENTRY_KIND: statement_only
% Imported from: batch16.tex; UnsolvedMath ID 5100002
\documentclass[11pt]{article}
\usepackage[margin=25mm]{geometry}
\usepackage{fontspec}
\setmainfont{Latin Modern Roman}
\usepackage{amsmath,amssymb,mathtools}
\usepackage{unicode-math}
\setmathfont{Latin Modern Math}
\usepackage{xurl,hyperref}
\hypersetup{colorlinks=true,urlcolor=blue}
\setcounter{secnumdepth}{0}
\setlength{\parindent}{0pt}
\setlength{\parskip}{5pt}
\setlength{\emergencystretch}{3em}
\newcommand{\sref}[2]{\hyperlink{source:#1:#2}{[S#2]}}
\newcommand{\cataloguescope}[1]{\paragraph{Recorded target.}#1}
\begin{document}
% UnsolvedMath ID 5100002; source catalogue code AMR-050-0002
\setcounter{section}{5100001}
\section{Elliptic-billiard invariant k\_\{108\}}\label{5100002}
{\small\sffamily UnsolvedMath ID 5100002 · Dynamical Systems}
\subsection{Definitions and mathematical statement}

\paragraph{Shared notation.}$\mathbb N=\{1,2,\ldots\}$, and $\mathbb Z,\mathbb Q,\mathbb R,\mathbb C$ denote the integers, rationals, reals and complex numbers. Any other conventions are specified locally.


Fix a noncircular ellipse
\[
E=\{(x,y)\in\mathbb R^2:x^2/a^2+y^2/b^2=1\},\qquad a>b>0,
\]
with center $O=(0,0)$ and foci $f_1=(\sqrt{a^2-b^2},0)$ and $f_2=-f_1$. Fix an inner confocal ellipse $C$, and an integer $N\geq3$, for which the pair $(E,C)$ admits a Poncelet family $\mathcal F$ of $N$-periodic billiard trajectories. A member $P=(P_1,\ldots,P_N)$ has vertices on $E$, each segment $P_iP_{i+1}$ is tangent to $C$, and incoming and outgoing segments satisfy the usual specular reflection law at $P_i$. Indices are cyclic: $P_{N+1}=P_1$.
For any ordered polygon $W=(W_1,\ldots,W_N)$ define its signed area by
\[
S(W)=\frac12\sum_{i=1}^N\det(W_i,W_{i+1}),\qquad
\det((x,y),(u,v))=xv-yu.
\]
Let $t_i$ be the tangent line to $E$ at $P_i$. The outer polygon $P'$ has vertices $P_i'=t_i\cap t_{i+1}$, and the inner polygon $P''$ has vertices $P_i''$, the contact points of $P_iP_{i+1}$ with $C$. Write $A=S(P)$, $A'=S(P')$ and $A''=S(P'')$. These areas retain their signs. Constancy below means constancy as $P$ varies in the fixed family $\mathcal F$; its value may depend on $E$, $C$ and $N$.
Let $\theta_i$ be the angle between $P_{i-1}-P_i$ and $P_{i+1}-P_i$, so that $\theta_i\in[0,\pi]$ and
\[
\cos\theta_i=\frac{(P_{i-1}-P_i)\cdot(P_{i+1}-P_i)}{|P_{i-1}-P_i|\,|P_{i+1}-P_i|}.
\]
Write $H(P)=\prod_{i=1}^N\sin(\theta_i/2)$. Here $|\cdot|$ and $\cdot$ denote Euclidean norm and inner product.
A displayed quotient is evaluated where its denominator is nonzero; its constancy means that its values at any two members where it is defined agree.
\paragraph{Statement recorded in the supplied source.}
For every such family with $N\equiv2\pmod4$, prove that $\displaystyle \frac{A'}{A\prod_{i=1}^N\sin(\theta_i/2)}$ is constant on $\mathcal F$. \sref{5100002}{2}
\subsection{Short English statement}
For every family with two more sides than a multiple of four, prove that the signed outer-to-orbit area ratio divided by the product of half-angle sines stays constant.
\subsection{Sources}
\begin{description}
\item[\hypertarget{source:5100002:1}{S1}] ULAM.AI and the UnsolvedMath Contributors, UnsolvedMath: A Curated Collection of Open Mathematics Problems, record 5100002 (AMR-050-0002). CC BY 4.0. \url{https://huggingface.co/datasets/ulamai/UnsolvedMath}.
\item[\hypertarget{source:5100002:2}{S2}] Dan Reznik, Ronaldo Garcia and Jair Koiller, Eighty New Invariants in the Elliptic Billiard (2020), arXiv version 11, section 3.1, Table 2, p. 5, invariant $k_{108}$; sections 1--2, pp. 1--3, and Appendix A, Table 12, p. 15. \url{https://arxiv.org/abs/2004.12497v11}.
\end{description}
\label{end:5100002}
\end{document}
